\chapter{Seguimiento y Desarrollo Grupal}

\section{Planificación del Trabajo}

\subsection*{Sprint 1 — Fundación: base de datos, autenticación y roles}

\textbf{Objetivo:}  
Setup del proyecto, schema de BD, autenticación, RLS y CRUD del menú.

\textbf{Historias de Usuario}

\begin{itemize}[leftmargin=1.5em]

\item \textbf{Setup e infraestructura}

\begin{itemize}
\item HU-01: Inicializar proyecto Expo con TypeScript y Expo Router
\item HU-02: Configurar proyecto Supabase y variables de entorno
\item HU-03: Crear migración de esquema inicial de la base de datos
\item HU-04: Crear constraint de orden activa única por mesa
\item HU-05: Crear trigger set\_closed\_at al entregar una orden
\end{itemize}

\item \textbf{Autenticación y roles}

\begin{itemize}
\item HU-06: Implementar pantalla de login para mesero y admin
\item HU-07: Proteger rutas según rol con guards de navegación
\item HU-08: Crear trigger sync\_role\_to\_auth para sincronizar rol
\item HU-09: Implementar cierre de sesión
\end{itemize}

\item \textbf{Row Level Security}

\begin{itemize}
\item HU-10: Implementar RLS para la tabla orders
\item HU-11: Implementar RLS para expenses, menu\_items y table\_calls
\end{itemize}

\item \textbf{Gestión de usuarios (admin)}

\begin{itemize}
\item HU-12: Listar usuarios del sistema (admin)
\item HU-13: Crear nuevo usuario mesero o admin (admin)
\item HU-14: Activar y desactivar usuario (admin)
\end{itemize}

\item \textbf{Gestión de mesas (admin)}

\begin{itemize}
\item HU-15: CRUD de mesas del restaurante (admin)
\item HU-16: Gestionar restaurante y su slug (admin)
\end{itemize}

\end{itemize}


\subsection*{Sprint 2 — Operación: menú, pedidos, pagos y gastos}

\textbf{Objetivo:}  
Registro de pedidos, vista de mesas en tiempo real y módulo financiero.

\textbf{Historias de Usuario}

\begin{itemize}[leftmargin=1.5em]

\item \textbf{Gestión de menú (admin)}

\begin{itemize}
\item HU-17: Listar ítems del menú con filtro por categoría (admin)
\item HU-18: Crear ítem del menú con imagen (admin)
\item HU-19: Editar precio y descripción de un ítem (admin)
\item HU-20: Activar o desactivar disponibilidad de un ítem (admin)
\item HU-21: Eliminar ítem del menú (admin)
\end{itemize}

\item \textbf{Toma de pedidos (mesero)}

\begin{itemize}
\item HU-22: Ver el menú disponible al tomar un pedido (mesero)
\item HU-23: Crear una nueva orden seleccionando mesa e ítems (mesero)
\item HU-24: Ver detalle de una orden activa (mesero)
\item HU-25: Agregar ítems a una orden existente (mesero)
\item HU-26: Marcar una orden como entregada (mesero)
\item HU-27: Cancelar una orden (mesero y admin)
\item HU-28: Seleccionar método de pago al crear la orden (mesero)
\end{itemize}

\item \textbf{Pagos (admin)}

\begin{itemize}
\item HU-29: Ver órdenes con pago QR pendiente de confirmación (admin)
\item HU-30: Confirmar pago QR de una orden (admin)
\item HU-31: Mostrar QR estático de pago al cliente (mesero)
\end{itemize}

\item \textbf{Gastos operativos (admin)}

\begin{itemize}
\item HU-32: Registrar un gasto operativo (admin)
\item HU-33: Listar gastos por rango de fechas (admin)
\item HU-34: Editar y eliminar un gasto (admin)
\end{itemize}

\end{itemize}


\subsection*{Sprint 3 — Tiempo real, carta digital y reportes}

\textbf{Objetivo:}  
Carta digital pública, sesión de cliente en mesa, pagos y despliegue final.

\textbf{Historias de Usuario}

\begin{itemize}[leftmargin=1.5em]

\item \textbf{Vista de mesas en tiempo real (mesero)}

\begin{itemize}
\item HU-35: Ver estado de todas las mesas asignadas (mesero)
\item HU-36: Actualizar estado de mesas en tiempo real sin recargar (mesero)
\item HU-37: Recibir notificación cuando un pedido está listo (mesero)
\item HU-38: Manejar reconexión automática del canal Realtime (mesero)
\end{itemize}

\item \textbf{Llamada al mesero (cliente → mesero)}

\begin{itemize}
\item HU-39: Mostrar botón "Llamar al mesero" en la sesión de cliente en mesa
\item HU-40: Recibir notificación de llamada de cliente en tiempo real (mesero)
\item HU-41: Marcar llamada de cliente como atendida (mesero)
\end{itemize}

\item \textbf{Carta digital pública (web)}

\begin{itemize}
\item HU-42: Publicar carta digital accesible por URL sin login
\item HU-43: Mostrar carta organizada por categorías con imágenes y precios
\item HU-44: Generar y mostrar QR del restaurante para carta pública (admin)
\end{itemize}

\item \textbf{Sesión de cliente en mesa (web)}

\begin{itemize}
\item HU-45: Iniciar sesión anónima al escanear el QR de la mesa
\item HU-46: Ver menú completo (incluyendo no disponibles) desde la sesión de mesa
\item HU-47: Generar QR individual por mesa (admin)
\end{itemize}

\item \textbf{Reportes financieros (admin)}

\begin{itemize}
\item HU-48: Crear función SQL get\_report para ingresos y egresos
\item HU-49: Crear función SQL get\_top\_items para ranking de platos
\item HU-50: Pantalla de reporte financiero por período (admin)
\item HU-51: Pantalla de ranking de platos más vendidos (admin)
\end{itemize}

\end{itemize}

\bigskip
\begin{center}
\textit{51 historias de usuario · 3 sprints · 6 semanas}
\end{center}


\section{Roles y Responsabilidades}

\begin{center}
\begin{tabular}{ll}
  \toprule
  \textbf{Integrante} & \textbf{Rol} \\
  \midrule
  Acebey Zalazar Jimmy Randall & Desarrollador \\
  Cáceres Véliz Erick Alvaro   & Desarrollador \\
  Revollo Bernal Amilcar Diego & Scrum Master \\
  Siles Espicho Adrian Eloy    & Desarrollador  \\
  Tinta Mamani Jairo Amado     & Desarrollador \\
  \bottomrule
\end{tabular}
\end{center}


\section{Control de Versiones}

El proyecto utiliza un sistema de control de versiones distribuido con
\textbf{Git}, gestionado en un único repositorio alojado en GitHub
(\texttt{gestion-comandas}). Este enfoque centralizado permite una
coordinación eficiente entre el desarrollo del frontend móvil, las rutas
web públicas y las migraciones de base de datos.

La estrategia de ramas adoptada es:

\begin{itemize}[leftmargin=1.5em]
  \item \textbf{\texttt{main}:} rama de producción; solo recibe merges mediante Pull Requests aprobados.
  \item \textbf{\texttt{feature/[nombre]}:} una rama por historia de usuario o tarea de sprint.
\end{itemize}


\section{Reporte de Avances}

\subsection{Indicadores de Progreso}

\begin{longtable}{>{\bfseries}p{2cm} p{9cm} c}
  \toprule
  \textbf{ID} & \textbf{Historia de Usuario} & \textbf{Completada} \\
  \midrule
  \endhead

  HU-01 & Inicializar proyecto Expo con TypeScript y Expo Router & \checkmark \\
  HU-02 & Configurar proyecto Supabase y variables de entorno & \checkmark \\
  HU-03 & Crear migración de esquema inicial de la base de datos & \checkmark \\
  HU-04 & Crear constraint de orden activa única por mesa & \checkmark \\
  HU-05 & Crear trigger set\_closed\_at al entregar una orden & \checkmark \\
  HU-06 & Implementar pantalla de login para mesero y admin & \checkmark \\
  HU-07 & Proteger rutas según rol con guards de navegación & \checkmark \\
  HU-08 & Crear trigger sync\_role\_to\_auth para sincronizar rol & \checkmark \\
  HU-09 & Implementar cierre de sesión & \checkmark \\
  HU-10 & Implementar RLS para la tabla orders & \checkmark \\
  HU-11 & Implementar RLS para expenses, menu\_items y table\_calls & \checkmark \\
  HU-12 & Listar usuarios del sistema (admin) & \checkmark \\
  HU-13 & Crear nuevo usuario mesero o admin (admin) & \checkmark \\
  HU-14 & Activar y desactivar usuario (admin) & \checkmark \\
  HU-15 & CRUD de mesas del restaurante (admin) & \checkmark \\
  HU-16 & Gestionar restaurante y su slug (admin) & \checkmark \\
  HU-17 & Listar ítems del menú con filtro por categoría (admin) & \checkmark \\
  HU-18 & Crear ítem del menú con imagen (admin) & \checkmark \\
  HU-19 & Editar precio y descripción de un ítem (admin) & \checkmark \\
  HU-20 & Activar o desactivar disponibilidad de un ítem (admin) & \checkmark \\
  HU-21 & Eliminar ítem del menú (admin) & \checkmark \\
  HU-22 & Ver el menú disponible al tomar un pedido (mesero) & \checkmark \\
  HU-23 & Crear una nueva orden seleccionando mesa e ítems (mesero) & \checkmark \\
  HU-24 & Ver detalle de una orden activa (mesero) & \checkmark \\
  HU-25 & Agregar ítems a una orden existente (mesero) & \checkmark \\
  HU-26 & Marcar una orden como entregada (mesero) & \checkmark \\
  HU-27 & Cancelar una orden (mesero y admin) & \checkmark \\
  HU-28 & Seleccionar método de pago al crear la orden (mesero) & \checkmark \\
  HU-29 & Ver órdenes con pago QR pendiente de confirmación (admin) & \checkmark \\
  HU-30 & Confirmar pago QR de una orden (admin) & \checkmark \\
  HU-31 & Mostrar QR estático de pago al cliente (mesero) & \checkmark \\
  HU-32 & Registrar un gasto operativo (admin) & \checkmark \\
  HU-33 & Listar gastos por rango de fechas (admin) & \checkmark \\
  HU-34 & Editar y eliminar un gasto (admin) & \checkmark \\
  HU-35 & Ver estado de todas las mesas asignadas (mesero) & \checkmark \\
  HU-36 & Actualizar estado de mesas en tiempo real sin recargar (mesero) & \checkmark \\
  HU-37 & Recibir notificación cuando un pedido está listo (mesero) & \checkmark \\
  HU-38 & Manejar reconexión automática del canal Realtime (mesero) & \checkmark \\
  HU-39 & Mostrar botón "Llamar al mesero" en la sesión de cliente en mesa & \checkmark \\
  HU-40 & Recibir notificación de llamada de cliente en tiempo real (mesero) & \checkmark \\
  HU-41 & Marcar llamada de cliente como atendida (mesero) & \checkmark \\
  HU-42 & Publicar carta digital accesible por URL sin login & \checkmark \\
  HU-43 & Mostrar carta organizada por categorías con imágenes y precios & \checkmark \\
  HU-44 & Generar y mostrar QR del restaurante para carta pública (admin) & \checkmark \\
  HU-45 & Iniciar sesión anónima al escanear el QR de la mesa & \checkmark \\
  HU-46 & Ver menú completo (incluyendo no disponibles) desde la sesión de mesa & \checkmark \\
  HU-47 & Generar QR individual por mesa (admin) & \checkmark \\
  HU-48 & Crear función SQL get\_report para ingresos y egresos & \checkmark \\
  HU-49 & Crear función SQL get\_top\_items para ranking de platos & \checkmark \\
  HU-50 & Pantalla de reporte financiero por período (admin) & \checkmark \\
  HU-51 & Pantalla de ranking de platos más vendidos (admin) & \checkmark \\

  \bottomrule
  \caption{Indicadores de progreso por historia de usuario}
\end{longtable}

\subsection{Dificultades Encontradas}

\begin{itemize}[leftmargin=1.5em]

\item \textbf{Pausa automática del free tier de Supabase:}
La base de datos se pausa tras períodos de inactividad,
por lo que se implementó un ping periódico para mantener
activo el servicio durante el desarrollo.

\item \textbf{Reconexión de canales Realtime:}
Las conexiones móviles inestables pueden interrumpir la
sincronización en tiempo real. Se implementó manejo de
eventos \texttt{CHANNEL\_ERROR} y \texttt{TIMED\_OUT}
con reintento automático.

\item \textbf{Persistencia de precios históricos:}
Para mantener la consistencia financiera, el precio del
producto se almacena como snapshot en el momento del pedido,
evitando inconsistencias cuando el admin modifica el menú.

\end{itemize}


\section{Comunicación}

El equipo empleó las siguientes herramientas de comunicación y colaboración:

\begin{itemize}[leftmargin=1.5em]
  \item \textbf{Notion:} planificación de tareas, backlog y documentación técnica.
  \item \textbf{GitHub:} control de versiones y revisión de código mediante pull requests.
  \item \textbf{Discord:} canal principal de comunicación durante el desarrollo.
  \item \textbf{WhatsApp:} coordinación rápida para decisiones puntuales.
\end{itemize}